% Part C; AI-produced by GPT-6 Astra (OpenAI), 19 November 2026.
% Intended for \input in the main paper. Requires amsmath, amssymb, amsthm,
% theorem environment, and natbib (or compatible \citep, \citet).
\section{Oriented maquettes and regular cell structures}
\label{sec:maquettes}

We distinguish the oriented set $M$, its quotient $P=M/\langle a\rangle$
by the fixed-point-free reversal $a$, and the order complex $\Delta(P)$.
The quotient map $q:M\to P$ forgets an orientation; it does not identify
points in a geometric cell. All order complexes have the weak topology.
The empty complex is denoted $S^{-1}$.

\subsection{A necessary audit of the literal axioms}

The transcription and our inspection of 156-8, scan~4, write Magor~$1'$(c) as
$X\mathrel{|\circ|}Y\Rightarrow X^-\mathrel{|\circ|}Y$;
scan~6 writes $X\mathrel{|\circ|}Y\Longleftrightarrow X\ne Y$.
These formulas cannot both hold on a nonempty oriented set.
We inspected both original scan images. The inequality glyph on scan~6
has no separate sign identifying orientation orbits. We retain the
literal deduction and the explicitly corrected interpretation separately.
The transcription also places a prime on the antecedent $Y$ of
Magor~$1'$(b), while the consequent uses $Y^-$. We read the same bound
variable on both sides. If the prime instead denotes an independently
quantified variable, identity refinement would force every $X$ to
equal every $Y$, again leaving only the empty case under free reversal.

\begin{theorem}[Proved: literal incompatibility]\label{SA21-C01}
Let $M$ have a fixed-point-free involution. If an irreflexive relation
$\perp\,\subseteq M\times M$ satisfies
$x\perp y\Rightarrow a(x)\perp y$, then it cannot equal inequality
unless $M=\varnothing$.
\end{theorem}
\begin{proof}
For $x\in M$, put $y=a(x)$. Inequality gives $x\perp a(x)$,
and reversal invariance gives $a(x)\perp a(x)$, contradicting
irreflexivity. The empty case satisfies every condition vacuously.
\end{proof}

The geometric correction used below is
$X\mathrel{|\circ|}Y\Longleftrightarrow q(X)\ne q(Y)$.
Changing Magor~$1'$(c) to simultaneous reversal is another consistent
correction; in the maquette case it gives the same incidence classification,
but it is a different general axiom. Likewise, the doubtful superscript
in Magor~2(b) is treated explicitly. The original image has a small
epsilon-like superscript on the source $X$, and no corresponding
superscript on the target $Y$. We do not identify this doubtful mark
with an unprinted target sign.

\begin{proposition}[Proved: the two readings of Magor 2(b)]
\label{SA21-C02}
Suppose interior refinement is identity and disjunction is corrected to
orbit inequality. If Magor~2(b) permits $X'\triangleleft X^\epsilon$
for either sign $\epsilon$ but concludes $X'\triangleleft Y$ with
$Y=X$, then Magor~1, $1'$, 5 and 6 permit precisely the graphs whose
edges are all loops; there are no cells of dimension at least two.
If the target is given the same sign, or the source superscript is
absent, Magor~2(b) is automatic and the general graph dictionary below
applies.
\end{proposition}
\begin{proof}
In the first reading, write $D_X$ for the set of immediate predecessors
of $X$. Taking $\epsilon=-1$, first at $X$ and then at $X^-$, gives
$D_X=D_{X^-}$. Simultaneous reversal then gives $a(D_X)=D_X$.
If $Z\triangleleft Y\triangleleft X$ has rank difference two,
both $Y$ and $Y^-$ lie between $Z$ and $X$, contradicting the uniqueness
in Magor~6(a). Any element of dimension at least two has an initial
length-two portion of a longest incidence chain; its last element has
dimension two. Thus such elements are impossible. For a one-cell,
Magor~6(b) and reversal-stability of $D_X$ force exactly the two signs
of a single vertex below it, so the associated edge is a loop.
Conversely arbitrary graphs consisting of loops and isolated vertices
satisfy all these axioms: rank-two conditions are empty and every
one-cell has precisely those two signed endpoints. In the same-sign
reading the asserted incidence is the assumed incidence because
$X=Y$; the reading without the source sign has the same property.
\end{proof}

Every geometric classification below uses this latter lifting convention,
orbit inequality, and the covering condition specified next. These are
stated corrections, not claims that the uncorrected formulas imply the
desired examples.
For Magor~6(a) we use the subject $Z$ in $\operatorname{codim}(Z,X)=2$,
as in the attached correction SA21-K8-09. If the doubtful small glyph
were instead an independently quantified $Y$, the universal closure of
that formula would allow no cell of dimension at least two: choose
$Y$ two ranks below $X$ on a longest chain and $Z$ one rank below $X$;
the demanded intermediate between $Z$ and $X$ would lengthen that
chain. Conversely the condition is vacuous on all graphs. Thus this
alternative reading cannot supply the intended higher-dimensional axiom.

We require an \emph{immediate} incidence to project to a cover relation
of the orbit poset. A merely acyclic generating relation need not have
this property: adding a redundant arrow across two ranks does not change
the generated order. The covering condition is therefore part of the
precise characterization below, unless already included in the definition
of immediate incidence.
The orbit relation is a partial order: any cycle of orbits can be lifted,
using simultaneous reversal, to a strict chain from $X$ to either $X$
or $X^-$. The first is impossible, and the second, followed by its
reversed-sign copy, is also a cycle. Thus the quotient order used below
is well defined.

\subsection{Connected boundaries and the low-dimensional dictionary}

\begin{theorem}[Proved: the connected-boundary axiom]\label{SA21-C03}
Let $P$ be any poset and $B_x=P_{<x}$. The following are equivalent:
\begin{enumerate}
\item $B_x$ is nonempty and connected for either Alexandrov topology;
\item its undirected comparability graph is connected;
\item $|\Delta(B_x)|$ is nonempty and connected.
\end{enumerate}
If principal ideals are finite, the comparability graph can be replaced
by the Hasse graph. Thus the exact additional incidence axiom for connected
cell boundaries is that any two proper faces of each cell of dimension
at least two can be joined by a finite incidence zigzag among its proper
faces. Dimensions zero and one must be excepted.
\end{theorem}
\begin{proof}
A comparability component is both upward and downward closed, hence
clopen for either Alexandrov topology. Conversely a clopen separation
cannot separate comparable elements: a lower-open set containing the
greater one contains the smaller one, and its lower-open complement
gives the reverse exclusion. The upper-open proof is dual. This proves
the first equivalence. Edges of the order complex are precisely comparable
pairs. Components of its edge graph partition all simplices and therefore
give disjoint open and closed subcomplexes in the weak topology; a finite
edge path is a topological path. This proves the second equivalence.
In a finite principal ideal every comparable pair is joined by a
saturated chain. Finally, a regular zero-cell has empty boundary and a
regular one-cell has boundary $S^0$, which is disconnected.
\end{proof}

\begin{theorem}[Proved under the stated corrected incidence conventions]
\label{SA21-C04}
In dimension at most one, oriented maquettes are precisely graphs given
by vertices $V$, oriented edges $E$, a free reversal $e\mapsto e^-$ on
$E$, and an origin map $s:E\to V$. In dimension at most two, the oriented
edges below each selected two-cell orientation form a disjoint union
of finite directed cycles and two-sided infinite directed paths, with
each incident vertex used once as origin and once as terminus.
If each cell has finitely many faces, only finite cycles occur.
\end{theorem}
\begin{proof}
Put $V=M_0^+$ and $E=M_1$. Magor~6(b) gives the unique origin
$s(e)$; applying it to $e^-$ and then reversing gives the unique
terminal vertex $t(e)=s(e^-)$. The incidences are exactly
$s(e)^+\triangleleft e$ and $t(e)^-\triangleleft e$, with their
simultaneous reversals. Conversely these formulas recover every incidence
from the graph data and verify the dimension-one axioms. This is Serre's
convention \citep[Chapter~I, \S2.1, p.~22]{Serre1977}; the attached
study attributes the identification to Michel Hua's project, which we
credit while giving the independent derivation.

Fix an orientation of a two-cell $f$ and let $A_f$ be the set of its
incident oriented edges. For each vertex below $f$, Magor~6(a) applied
to its two orientations gives exactly one outgoing and one incoming
edge in $A_f$. Thus following outgoing edges defines a permutation of
the incident vertex set. Each orbit of a permutation is either a finite
cycle or an infinite orbit indexed bijectively by $\mathbb Z$.
Conversely such a union has exactly the two required signed two-step
paths at every vertex and so verifies Magor~6(a). With a finite face set
the infinite orbits are impossible.
\end{proof}

This corrects the unrestricted finite-cycle formulation of R9. A
two-sided path with vertices $v_i$, edges $e_i:v_i\to v_{i+1}$
($i\in\mathbb Z$), and a formal two-cell having every $e_i$ in its
positive boundary is a connected infinite counterexample. Its
incidence chains have length at most two, but its boundary is a line,
not a circle; finite chain length does not imply finite cell closure.

\begin{counterexample}[Refuted: excluding loops alone is insufficient]
\label{SA21-C05}
Even in finite dimension two, connected boundaries and the exclusion
of loops do not make an oriented maquette regular. With finite faces
per cell, the exact dimension-at-most-two regular class additionally
excludes using both orientations of an underlying edge in one two-cell
boundary.
\end{counterexample}
\begin{proof}
Take two vertices $v_0,v_1$, one non-loop edge $e$ between them,
and one formal two-cell $f$ with both $e^+$ and $e^-$ incident to
$f^+$. Include simultaneous reversals. More explicitly the signed
vertex incidences are
\[
 v_0^+\triangleleft e^+,\quad v_1^-\triangleleft e^+,
 \qquad
 v_0^-\triangleleft e^-,\quad v_1^+\triangleleft e^-.
\]
Each sign of $e$ is incident to each sign of $f$. For either sign of
either vertex there is exactly one two-step incidence path to each
sign of $f$, and likewise one from the opposite vertex orientation.
Thus Magor~6 holds. The unoriented boundary poset of $f$ has two
vertices below a single edge; its order complex is an interval and
cannot decompose the boundary circle of a regular two-cell. This
eight-element signed structure is checked exhaustively by the supplied
program.

For the corrected positive statement, a finite connected directed
boundary with one incoming and one outgoing edge per vertex is a cycle.
When it has no loop and uses no underlying edge twice, it embeds as a
circle in the underlying graph. Attach a disk along that homeomorphism.
Doing this for every two-cell yields a CW complex with all closed cells
embedded. Conversely regularity gives exactly these boundary embeddings
and finite face sets. The number of cells and the valence at a vertex
may be infinite.
\end{proof}

\subsection{The three-dimensional obstruction}

\begin{counterexample}[Refuted: Lead L4]\label{SA21-C06}
Finite signed diamonds, connected positive-dimensional cell boundaries,
and connected vertex links in the order complex do not characterize
regular CW face posets. There is a counterexample with $55$ unoriented
cells, of maximal dimension three.
\end{counterexample}
\begin{proof}
Let $V=(\mathbb Z/3\mathbb Z)^2$. For each $(i,j)$ take the oriented
triangles
\[
 [(i,j),(i+1,j),(i+1,j+1)],\qquad
 [(i,j),(i+1,j+1),(i,j+1)].
\]
Let $T$ be the poset of their nonempty faces, and adjoin a greatest
element $z$ of dimension three. These triangles are the standard
periodic $3\times3$ triangulation of a torus. Orient edges by increasing
vertex label. Use alternating simplex incidence signs below triangles,
and for $[f:z]$ use the sign of the displayed orientation of $f$ relative
to its increasing order. Double every cell by signs and set
\[
 (a,s)\triangleleft(b,t)
 \quad\Longleftrightarrow\quad
 a\text{ is a facet of }b\quad\text{and}\quad st[a:b]=1.
\]
The rank-two intervals, after adding an empty face, are the two
endpoints of an edge, the two edges of a triangle through a vertex,
and the two triangles incident with an edge of the torus. Their sign
products are opposite. There are $9$ vertices, $27$ edges, $18$
triangles, and the formal cell $z$. Every surface vertex has a six-cycle
link; all triangles and the torus have connected boundaries of the
dimensions required. In $\Delta(T\cup\{z\})$ the link of $z$ is the
subdivided torus. At any proper face the upper interval has maximum
$z$, so the order-complex link has a nonempty cone factor and is
connected. This verifies the specified link condition too.

The top boundary has rational Betti numbers $(1,2,1)$ and Euler
characteristic $9-27+18=0$. A regular three-cell has boundary $S^2$,
whose first homology is zero and whose Euler characteristic is $2$.
Hence this poset cannot be a regular CW face poset. The program constructs
all incidence relations, verifies all $108$ rank-two intervals and
$324$ oriented instances of Magor~6(a), and calculates the rational
boundary-matrix ranks $8,17$. No classification of surfaces is needed
for the homological obstruction.
\end{proof}

If ``every cell link is connected'' is interpreted to include a
zero-dimensional spherical link, the proposed condition is not even
necessary for regular CW complexes: such a link is $S^0$. The preceding
counterexample addresses the usual nontrivial interpretations that
exclude that unavoidable exception. Higher-dimensional link
connectedness cannot replace spherical lower intervals.

\begin{counterexample}[Refuted: connected global links are necessary]
\label{SA21-C07}
Connectedness of positive-dimensional global links is not necessary for
a poset to be a regular CW face poset.
\end{counterexample}
\begin{proof}
Take the abstract simplicial complex with facets $\{0,1,2\}$ and
$\{0,3,4\}$. This is a regular CW complex with five vertices, six edges
and two triangles. The link of the common vertex $0$ consists of the
two disjoint edges $\{1,2\}$ and $\{3,4\}$, so it is one-dimensional
and disconnected. Its link in the barycentric order complex likewise
has two components. The finite checker verifies these assertions.
The lower interval below a cell, an upper interval of cofaces, and a
global link must therefore not be treated as interchangeable.
\end{proof}

\subsection{An exact characterization in all dimensions}

For the signed condition, choose one lift of each orientation orbit,
using the distinguished positive lift at vertices. For a cover $y<x$,
write $\epsilon(y,x)=+1$ when these chosen lifts are incident and
$\epsilon(y,x)=-1$ when the opposite lower lift is incident to the
chosen upper lift. This notation is defined by the non-repetition
condition below. Changing a chosen lift changes the adjoining signs
but preserves the signed-diamond equations.

\begin{theorem}[Proved; the unoriented classification is known]
\label{SA21-C08}
Let $M$ carry free reversal and an immediate incidence relation invariant
under simultaneous reversal. Suppose the induced relation on orientation
orbits generates a poset $P$. These data arise from a regular CW complex
with oriented cells and distinguished positive zero-cell orientations
if and only if the following conditions hold:
\begin{enumerate}
\item Every principal ideal of $P$ is finite.
\item Immediate incidences project exactly to covers of $P$; for each
cover $u<v$ and each orientation of $v$, exactly one orientation of $u$
is incident to it.
\item For every $x\in P$, $|\Delta(P_{<x})|$ is homeomorphic to
$S^{d(x)-1}$ for some integer $d(x)\ge0$.
\item The products of the two incidence signs around every rank-two
diamond are opposite; at an edge the two vertex coefficients have
opposite signs relative to the distinguished vertex orientations.
\end{enumerate}
The spherical condition implies that the poset is graded and that
$d(x)$ is its dimension rank. The theorem
allows infinitely many cells, unbounded dimension, and failure of local
finiteness. The empty case gives the empty CW complex.
\end{theorem}
\begin{proof}
For the unoriented statement apply \citet[Definition~2.1, p.~8;
Proposition~3.1, p.~11]{Bjorner1984} to $P$ with an adjoined minimum.
Its hypotheses are precisely the spherical lower-interval conditions
with ranks shifted by one. The associated realization has closed cell
$|\Delta(P_{\le x})|$ and open cell
\[
 |\Delta(P_{\le x})|\setminus |\Delta(P_{<x})|.
\]
The closure is the cone on the stipulated sphere. Finite ideals ensure
closure-finiteness. To make the construction explicit, every point of
the order complex has a finite supporting chain and hence a unique
largest supporting element $x$; it lies in exactly the open cell
displayed above. Each such open cone is an open ball, and its boundary
is the union of the cells indexed by $y<x$. The weak topology for these
cells equals the simplicial weak topology: each closed cell is a finite
subcomplex and each simplex lies in the closed cell of its largest
vertex. Consequently closedness can be tested by either family of
closed pieces. This is a regular CW structure even when upper intervals
or the global dimension are infinite. The map taking
$x$ to this open cell is an isomorphism of face posets. We use a
homeomorphism of cellular realizations here, not a homeomorphism to
the Alexandrov space $P$.

We must additionally show that the given signs, rather than merely
some signs, are realized. Choose orientations by dimension. The
zero-cell orientations are prescribed, and the edge condition permits
the choices in dimension one. Suppose all lower-dimensional orientations
have been chosen. For an $n$-cell $x$ the prescribed finite boundary chain
is
\[
 c_x=\sum_{y\lessdot x}\epsilon(y,x)y.
\]
Each coefficient of its boundary is the sum of the two path products
in a signed diamond and hence vanishes. Thus $c_x$ is a top-dimensional
cellular cycle in the boundary $(n-1)$-sphere. This cycle group is
$\mathbb Z$: there are no $n$-cells in the boundary, and its top homology
is infinite cyclic. A fundamental cycle has coefficient $+1$ or $-1$
at every top cell. Since the same is true of $c_x$, it equals one of
the two fundamental generators. Choose the orientation of $x$ giving
that generator. This is independent for each $n$-cell, and induction
realizes every prescribed sign.

Conversely, regular characteristic maps give finite cell closures and
spherical boundaries. Applying the subdivision construction inductively
on cell dimension identifies the order complex of each proper-face
poset with that boundary sphere. There is a unique signed incidence
along each facet,
and the cellular identity $\partial^2=0$. Rank-two intervals in a
regular CW face poset have exactly two intermediate cells, so this
identity is the required opposite-product equation. This proves all
conditions necessary.
\end{proof}

This is an object classification with an explicit comparison: the
signed double of the face-poset isomorphism is an involution- and
incidence-preserving bijection from $M$ to the oriented cells. No
equivalence of categories with unspecified morphisms is asserted.
The proof settles the corrected version of C.1 in every dimension.

\subsection{Polyhedral and simplicial conditions}

\begin{theorem}[Known, with the exact application proved]
\label{SA21-C09}
For a poset $P$ of nonempty cells, let $P^0=P\cup\{0\}$ with a new
minimum. The following are equivalent:
\begin{enumerate}
\item $P$ is the nonempty face poset of an abstract simplicial complex;
\item every interval $[0,x]$ is a finite Boolean lattice and $P^0$ is a
meet-semilattice;
\item every interval $[0,x]$ is finite Boolean and the map
$x\mapsto V_x$, where $V_x$ is the set of atoms below $x$, is injective.
\end{enumerate}
The equivalence of the first two conditions is stated by
\citet[\S2.3, pp.~8--9]{Bjorner1984}.
\end{theorem}
\begin{proof}
We verify the extra formulation and the map needed in this application.
Under local Booleanity, every element in a principal ideal is determined
by its atoms. If global meets exist and $V_x=V_y$, then $x\wedge y$
has this same atom set; therefore $x=x\wedge y=y$. Conversely, with
injectivity every subset $S\subseteq V_x$ is the atom set of the unique
element corresponding to it in $[0,x]$. Hence the image is a hereditary
family of finite nonempty subsets. If $V_x\subseteq V_y$, the element
below $y$ having atom set $V_x$ must be $x$, so the map reflects order.
It is consequently a poset isomorphism onto an abstract simplicial face
poset, and intersections supply the global meets after adjoining $0$.
The empty case is immediate.
\end{proof}

Local Booleanity alone gives simplicial posets. For example, two parallel
edges with the same two endpoints form a regular circle whose principal
augmented ideals are Boolean, but whose edge cells have the same atom
set. This explicitly checked four-cell example explains why a global
intersection condition is also needed.

\begin{theorem}[Refuted and corrected: At pol 1 plus At simp 1]
\label{SA21-C10}
For maquettes, At pol~1 (156-7/14) says exactly that $P^0$ is a
meet-semilattice. Within the regular CW class, At simp~1 (156-7/15)
is automatic. Their conjunction does not characterize simplices or
simplicial complexes. The square is a finite counterexample. An exact
replacement is At pol~1 together with finite Boolean principal ideals.
Equivalently, one can require finite distributive principal augmented
lattices together with At simp~1, retaining At pol~1.
\end{theorem}
\begin{proof}
In a maquette all pairs are compatible, so At pol~1 gives a greatest
common lower bound whenever there is any common lower bound. The
adjoined empty element handles the remaining pairs. For a regular cell,
a nonempty proper-face ideal with only one maximum would have a conical
order complex, contradicting its spherical realization. Thus At simp~1
holds. The square has four vertices, four edges and one top face; pairwise
face intersections are empty or a face and each non-minimal face has
at least two maximal proper faces. However, its principal lattice has
$10$ elements, whereas a Boolean lattice with its four atoms has $16$.
This proves the refutation and rechecks SA21-K7-03.
Even the existence of a regular cell structure does not follow from
these two axioms: three incomparable vertices below one top element
satisfy both, but the top has three boundary points rather than the
two required for a regular one-cell. This four-cell example is also
checked by the program.

For the final equivalent repair, in a finite distributive lattice the
condition on lower covers means that every nonzero join-irreducible is
an atom. Every element is a join of join-irreducibles, so the lattice
is atomistic. The map to subsets of atoms is then a Boolean isomorphism:
distributivity makes distinct atoms disjoint, and every subset has its
join. Conversely finite Boolean lattices satisfy both requirements.
The preceding theorem now applies to the whole maquette. This also
independently checks SA21-R7-06.
\end{proof}

Within regular CW structures, At pol~1 means that two cell closures
intersect in either the empty set or one closed cell. It asserts no
realizability by convex polytopes. In a general atelier, the same
condition is restricted to compatible pairs, as the manuscript says.
The original definition immediately before the erroneous shortcut
already requires each principal ideal of nonempty faces to be a finite
nonempty-subset lattice (156-7/14, lines~139--143 of
\texttt{J156-7-01\_transcription.tex}). That definition, together with
At pol~1, is correct. It is the asserted replacement by the predecessor
count, and its subsequent use as At simp~1, that the square refutes.

\subsection{Flags, interval graphs, and the comparison maps}

\begin{theorem}[Proved: the two flag constructions]\label{SA21-C11}
For any poset $L$, the finite nonempty flags ordered by inclusion form
the simplicial face poset of $\Delta(L)$. By contrast, the construction
of 156-9/5--7 has vertex cells $a\in L$ and interval cells $I_{a,b}$
for $a<b$; its incidence maquette is the one-dimensional simplicial
complex given by the comparability graph of $L$. Its finite prefigures
are positional flags, and its finite figures are exactly those flags
closed under adjoining the two endpoints of every interval they contain.
\end{theorem}
\begin{proof}
Subsets of a finite flag are flags. Its principal augmented ideal is
therefore Boolean, and intersections supply global meets. Theorem
\ref{SA21-C09} gives the first assertion in every dimension.
In the second construction, incidence is defined by
$a,b<_{\rm face}I_{a,b}$, so there is one graph edge for each comparable
pair and there are no higher-dimensional cells. The positional order
is a different relation, with
$a<_{\rm pos}I_{a,b}<_{\rm pos}b$.
The manuscript's exact corollary on scan~7 concerns prefigures. Its
closure map is
\[
 \Phi\longmapsto\overline\Phi
 =\Phi\cup\bigcup_{I_{a,b}\in\Phi}\{a,b\}.
\]
If a member precedes $I_{a,b}$ positionally, its last endpoint is at
most $a$; if it follows, its first endpoint is at least $b$. Thus
inserting the endpoints preserves the flag property. An incidence-closed
flag is a figure, and a figure has this closure property. In particular
$\{I_{a,b}\}$ is a flag and prefigure but is not a figure. This qualifies
the unclosed wording in R13.
\end{proof}

For a regular CW poset the flag construction is its barycentric
subdivision, with a homeomorphism to the same Hausdorff realization.

\begin{theorem}[Known; the application and comparison are proved]
\label{SA21-C12}
For an arbitrary Alexandrov poset, the map
\[
 \mu_L:|\Delta(L)|\longrightarrow L,
 \qquad \mu_L(p)=\min\operatorname{supp}(p),
\]
for the lower-open convention, is a weak homotopy equivalence:
\citet[Theorem~2, p.~466; construction in \S5, p.~469]{McCord1966}.
This includes infinite posets and complexes with the weak topology;
see also \citet[Theorem~1.4.6 and Remark~1.4.16]{Barmak2011}.
It is generally neither a homeomorphism nor a homotopy equivalence.
\end{theorem}
\begin{proof}

For completeness, we verify the hypotheses of McCord's local theorem
in this application. Put $U_x=L_{\le x}$. The preimage
$\mu_L^{-1}(U_x)$ consists of points having a positive coordinate at
some vertex of $U_x$, and is the complement of the full subcomplex
on $L\setminus U_x$; it is open. Renormalizing the coordinates on
$U_x$ gives a deformation retraction onto $|\Delta(U_x)|$, a cone
with apex $x$. Each $U_x$ is contractible as an Alexandrov space and
these sets form a basis. Thus all restricted maps are weak equivalences,
and McCord's Theorem~6, p.~467, applies. This argument needs no local
finiteness. To see concretely why a homotopy equivalence need not result,
take the four-element poset with two incomparable minima below two
incomparable maxima. Its order complex is a circle. Every continuous
map from this connected finite poset to a Hausdorff space is constant:
comparable points must have the same image, and the comparability graph
is connected. A homotopy inverse into the circle would therefore make
the circle's identity null-homotopic, an impossibility. The checker
verifies that the four-element order complex is the four-cycle.
\end{proof}

The subdivision homeomorphism for a regular CW realization and the
weak equivalence to an Alexandrov poset have different types. They do
not identify the Alexandrov and Hausdorff topologies.
